\Chapter{56}

\Verse{1}
{
    \zA{话说平儿陪着凤姐吃了饭，伏侍盥漱毕，方往探春处来，只见院中寂静，只有 丫鬟婆子一个个都站在窗外听候。 平儿进入厅中，他姐妹姑嫂三人正商议些家务，
        说的便是年内赖大家请吃酒，他家花园中事故。 见他来了，探春便命他脚踏上坐了， 因说道：“我想的事，不为别的，只想着我们一月所用的头油脂粉又是二两的事。
        我想咱们一月已有了二两月银，丫头们又另有月钱，可不是又同刚才学里的八两一 样重重叠叠？}
}
{
    \eA{But let us pick up the clue of our story. P'ing Erh bore lady Feng company
        during her meal; then attending to her, while she rinsed her mouth and
        washed her hands, she betook herself eventually to T'an Ch'un's quarters,
        where she discovered the courtyard in perfect stillness.}
}
{
}

\Verse{2}
{
    \zA{这事虽小，钱有限，看起来也不妥当，你奶奶怎么就没想到这个呢？” 平儿笑道：“这有个原故：姑娘们所用的这些东西，自然该有分例，每月每处买办
        买了，令女人们交送我们收管，不过预备姑娘们使用就罢了，没有个我们天天各人 拿着钱，找人买这些去的。 所以外头买办总领了去，按月使女人按房交给我们。}
}
{
    \eA{Not a soul was about beyond several maids, matrons and close attendants of
        the inner rooms, who stood outside the windows on the alert to obey any
        calls. P'ing Erh stepped into the hall. The two cousins and their
        sister-in-law were all three engaged in discussing some domestic affairs.}
}
{
}

\Verse{3}
{
    \zA{至 于姑娘们每月的这二两，原不是为买这些的，为的是一时当家的奶奶太太，或不在 家，或不得闲，姑娘们偶然要个钱使，省得找人去：这不过是恐怕姑娘们受委屈意
        思。 如今我冷眼看着，各屋里我们的姐妹都是现拿钱买这些东西的，竟有了一半子。 我就疑惑不是买办脱了空，就是买的不是正经货。” 探春李纨都笑道：“你也留心
        看出来了。 脱空是没有的，只是迟些日子，催急了，不知那里弄些来，不过是个名 儿。}
}
{
    \eA{They were talking about the feast, to which they had been invited during the
        new year festivities by Lai Ta's wife, and various details in connection
        with the garden she had in her place. But as soon as she (P'ing Erh)
        appeared on the scene, T'an Ch'un desired her to seat herself on her
        footstool. "What was exercising my mind," she thereupon observed, "confines
        itself to this.}
}
{
}

\Verse{4}
{
    \zA{其实使不得，依然还得现买，就用二两银子，另叫别人的奶妈子的弟兄儿子买 来方才使得。 要使官中的人去，依然是那一样的，不知他们是什么法子？” 平儿便
        笑道：“买办买的是那东西，别人买了好的来，买办的也不依他，又说他使坏心， 要夺他的买办。 所以他们宁可得罪了里头，不肯得罪了外头办事的。 要是姑娘们使
        了奶妈子们，他们也就不敢说闲话了。” 探春道：“因此我心里不自在，饶费了两起钱，东西又白丢一半。 不如竟把买 办的这一项每月蠲了为是。}
}
{
    \eA{I was computing that the head-oil, and rouge and powder, we use during the
        course of a month, are also a matter of a couple of taels;}
}
{
}

\Verse{5}
{
    \zA{此是第一件事。 第二件，年里往赖大家去。 你也去的： 你看他那小园子比咱们这个如何？” 平儿笑道：“还没有咱们这一半大，树木花草 也少多着呢。”
        探春道：“我因和他们家的女孩儿说闲话儿，他说这园子除他们带 的花儿，吃的笋菜鱼虾，一年还有人包了去，年终足有二百两银子剩。 从那日，我
        才知道一个破荷叶、一根枯草根子，都是值钱的。” 宝钗笑道：“真真膏粱纨？}
}
{
    \eA{and I was thinking that what with the sum of two taels, already allotted us
        every month, and the extra monthly amount given as well to the maids,
        allowances are, with the addition again of that of eight taels for school
        expenses, we recently spoke about, piled to be sure one upon another. The
        thing is, it's true, a mere trifle, and the amount only a bagatelle, but it
        doesn't seem to be quite proper.}
}
{
}

\Verse{6}
{
    \zA{之 谈!你们虽是千金，原不知道这些事，但只你们也都念过书，识过字的，竟没看见 过朱夫子有一篇‘不自弃’的文么？” 探春笑道：“虽也看过，不过是勉人自励，
        虚比浮词，那里真是有的？” 宝钗道：“朱子都行了虚比浮词了？ 那句句都是有的。 你才办了两天事，就利欲熏心，把朱子都看虚浮了。
        你再出去，见了那些利弊大事， 越发连孔子也都看虚了呢！” 探春笑道：“你这样一个通人，竟没看见姬子书？}
}
{
    \eA{But how is it that your mistress didn't take this into account?" P'ing Erh
        smiled. "There's a why and a wherefore," she answered. "All the things
        required by you, young ladies, must absolutely be subject to a fixed rule;
        for the different compradores have to lay in a stock of each every month;
        and to send them to us by the maids to take charge of; but purely and simply
        to keep in readiness for you to use.}
}
{
}

\Verse{7}
{
    \zA{当 日姬子有云：‘登利禄之场，处运筹之界者，穷尧舜之词，背孔孟之道。’” 宝钗 笑道：“底下一句呢？” 探春笑道：“如今断章取意；
        念出底下一句，我自己骂我 自己不成？” 宝钗道：“天下没有不可用的东西，既可用，便值钱。 难为你是个聪 明人，这大节目正事竟没经历。”
        李纨笑道：“叫人家来了，又不说正事，你们且 对讲学问！” 宝钗道：“学问中便是正事。 若不拿学问提着，便都流入市俗去了。” 三人取笑了一回，便仍谈正事。}
}
{
    \eA{No such thing could ever be tolerated as that each of us should have to get
        money every day and try and hunt up some one to go and buy these articles
        for us! That's how it is that the compradores outside receive a lump sum,
        and that they send us, month by month, by the female servants the supplies
        allotted for the different rooms.}
}
{
}

\Verse{8}
{
    \zA{探春又接说道：“咱们这个园子，只算比他们 的多一半，加一倍算起来，一年就有四百银子的利息。 若此时也出脱生发银子，自 然小器，不是咱们这样人家的事。
        若派出两个一定的人来，既有许多值钱的东西， 任人作践了，也似乎暴殄天物。 不如在园子里所有的老妈妈中，拣出几个老成本分、 能知园圃的，派他们收拾料理。
        也不必要他们交租纳税，只问他们一年可以孝敬些 什么。 一则园子有专定之人修理花木，自然一年好似一年了，也不用临时忙乱；}
}
{
    \eA{As regards the two taels monthly allowed you, young ladies, they were not
        originally intended that you should purchase any such articles with, but
        that you should, if at any time the ladies in charge of the household
        affairs happened to be away from home or to have no leisure, be saved the
        trouble of having to go in search of the proper persons, in the event of
        your suddenly finding yourselves in need of money.}
}
{
}

\Verse{9}
{
    \zA{二 则也不致作践，白辜负了东西； 三则老妈妈们也可借此小补，不枉成年家在园中辛 苦； 四则也可省了这些花儿匠、山子匠并打扫人等的工费。
        将此有馀，以补不足， 未为不可。” 宝钗正在地下看壁上的字画，听如此说，便点头笑道：“善哉!‘三 年之内，无饥馑矣。’”
        李纨道：“好主意!果然这么行，太太必喜欢。 省钱事小， 园子有人打扫，专司其职，又许他去卖钱，使之以权，动之以利，再无不尽职的了。”
        平儿道：“这件事须得姑娘说出来。}
}
{
    \eA{This was done simply because it was feared that you would be subjected to
        inconvenience. But an unprejudiced glance about me now shows me that at
        least half of our young mistresses in the various quarters invariably
        purchase these things with ready money of their own; so I can't help
        suspecting that, if it isn't a question of the compradores shirking their
        duties, it must be that what they buy is all mere rubbish."}
}
{
}

\Verse{10}
{
    \zA{我们奶奶虽有此心，未必好出口。 此刻姑 娘们在园里住着，不能多弄些玩意儿陪衬，反叫人去监管修理，图省钱，这话断不 好出口。”
        宝钗忙走过来，摸着他的脸笑道：“你张开嘴，我瞧瞧你的牙齿舌头是 什么做的？ 从早起来到这会子，你说了这些话，一套一个样子：也不奉承三姑娘，
        也不说你们奶奶才短想不到； 三姑娘说一套话出来，你就有一套话回奉，总是三姑 娘想得到的，你们奶奶也想到了，只是必有个不可办的原故。}
}
{
    \eA{T'an Ch'un and Li Wan laughed. "You must have kept a sharp lookout to have
        managed to detect these things!" they said. "But as for shirking the
        purchases, they don't actually do so. It's simply that they're behind time
        by a good number of days. Yet when one puts on the screw with them, they get
        some articles from somewhere or other, who knows where? These are however
        only a sham;}
}
{
}

\Verse{11}
{
    \zA{这会子又是因姑娘们 住的园子，不好因省钱令人去监管。 你们想想这话，要果真交给人弄钱去的，那人
        自然是一枝花也不许掐，一个果子也不许动了，姑娘们分中自然是不敢讲究，天天 和小姑娘们就吵不清。 他这远愁近虑，不亢不卑，他们奶奶就不是和咱们好，听他
        这一番话，也必要自愧的变好了。” 探春笑道：“我早起一肚子气，听他来了，忽 然想起他主子来：素日当家，使出来的好撒野的人!我见了他更生气了。 谁知他来
        了，避猫鼠儿似的，站了半日，怪可怜的。}
}
{
    \eA{for, in reality, they aren't fit for use. But as they're now as ever
        obtained with cash down, a couple of taels could very well be given to the
        brothers or sons of some of the other people's nurses to purchase them with.
        They'll then be good for something! Were we however to employ any of the
        public domestics in the establishment, the things will be just as bad as
        ever.}
}
{
}

\Verse{12}
{
    \zA{接着又说了那些话，不说他主子待我好， 倒说‘不枉姑娘待我们奶奶素日的情意了’，这一句话，不但没了气，我倒愧了， 又伤起心来。
        我细想：我一个女孩儿家，自己还闹得没人疼没人顾的，我那里还有 好处去待人？” 口内说到这里，不免又流下泪来。 李纨等见他说得恳切，又想他素
        日赵姨娘每生诽谤，在王夫人跟前，亦为赵姨娘所累，也都不免流下泪来，都忙劝 他：“趁今日清净，大家商议两件兴利剔弊的事情，也不枉太太委托一场。 又提这
        没要紧的事做什么。”}
}
{
    \eA{I wonder how they do manage to get such utter rot as they do?" "The
        purchases of the compradores may be what they are," P'ing Erh smiled; "but
        were anyone else to buy any better articles, the compradores themselves
        won't ever forgive them. Besides other things, they'll aver that they
        harbour evil designs, and that they wish to deprive them of their post.}
}
{
}

\Verse{13}
{
    \zA{平儿忙道：“我已明白了。 姑娘说谁好，竟一派人就完了。” 探春道：“虽如此说，也须得回你奶奶一声儿。 我们这里搜剔小利，已经不当。 皆
        因你奶奶是个明白人，我才这样行； 若是糊涂多歪多妒的，我也不肯，倒像抓他的 乖的似的。 岂可不商议了行呢？” 平儿笑道：“这么着，我去告诉一声儿。”
        说着 去了； 半日方回来，笑道：“我说是白走一趟。 这样好事，奶奶岂有不依的！”}
}
{
    \eA{That's how it comes about that the servants would much rather give offence
        to you all inside, (by getting inferior things), and that they have no
        desire to hurt the feelings of the managers outside, (by purchasing anything
        of superior quality). But if you, young ladies, requisition the services of
        the nurses, these men won't have the arrogance to make any nonsensical
        remarks."}
}
{
}

\Verse{14}
{
    \zA{探春听了，便和李纨命人将园中所有婆子的名单要来，大家参度，大概定了几 个人。 又将他们一齐传来，李纨大概告诉给他们。 众人听了，无不愿意。 也有说：
        “那片竹子单交给我，一年工夫，明年又是一片。 除了家里吃的笋，一年还可交些 钱粮。” 这一个说：“那一片稻地交给我，一年这些玩的大小雀鸟的粮食，不必动
        官中钱粮，我还可以交钱粮。” 探春才要说话，人回：“大夫来了，进园瞧史姑娘 去。” 众婆子只得去领大夫。}
}
{
    \eA{"This accounts for the unhappy state my heart is in," T'an Ch'un observed.
        "But as we're called upon to squander money right and left, and as the
        things purchased are half of them uselessly thrown away, wouldn't it, after
        all, be better for us to eliminate this monthly allowance to the
        compradores? This is the first thing. The next I'd like to ask you is this.}
}
{
}

\Verse{15}
{
    \zA{平儿忙说：“单你们，有一百也不成个体统。 难道没 有两个管事的头脑儿带进大夫来？” 回事的那人说：“有吴大娘和单大娘，他两个 在西南角上聚锦门等着呢。”
        平儿听说，方罢了。 众婆子去后，探春问宝钗：“如何？” 宝钗笑答道：“幸于始者怠于终，善其 辞者嗜其利。”
        探春听了，点头称赞，便向册上指出几个来与他三人看。 平儿忙去 取笔砚来。}
}
{
    \eA{When they went, during the new year festivities, to Lai Ta's house, you also
        went with them; and what do think of that small garden as compared with this
        of ours?" "It isn't half as big as ours," P'ing Erh laughingly explained.
        "The trees and plants are likewise fewer by a good deal."}
}
{
}

\Verse{16}
{
    \zA{他三人说道：“这一个老祝妈，是个妥当的，况他老头子和他儿子，代 代都是管打扫竹子，如今竟把这所有的竹子交与他。 这一个老田妈本是种庄稼的，
        稻香村一带，凡有菜蔬稻稗之类，虽是玩意儿，不必认真大治大耕，也须得他去再 细细按时加些植养，岂不更好？” 探春又笑道：“可惜蘅芜院和怡红院这两处大地
        方，竟没有出息之物。” 李纨忙笑道：“蘅芜院里更利害，如今香料铺并大市大庙 卖的各处香料香草儿，都不是这些东西？ 算起来，比别的利息更大。}
}
{
    \eA{"When I was having a chat with their daughter," T'an Ch'un proceeded, "she
        said that, besides the flowers they wear, and the bamboo shoots, vegetables,
        fish and shrimps they eat from this garden of theirs, there's still enough
        every year for people to take over under contract, and that at the close of
        each year there's a surplus in full of two hundred taels.}
}
{
}

\Verse{17}
{
    \zA{怡红院别说别 的，单只说春夏两季的玫瑰花，共下多少花朵儿？ 还有一带篱笆上的蔷薇、月季、
        宝相、金银花、藤花，这几色草花，干了卖到茶叶铺药铺去，也值好些钱。” 探春 笑着点头儿，又道：“只是弄香草没有在行的人。” 平儿忙笑道：“跟宝姑娘的莺
        儿他妈，就是会弄这个的。 上回他还采了些晒干了，编成花篮葫芦给我玩呢。 姑娘 倒忘了么？” 宝钗笑道：“我才赞你，你倒来捉弄我了。”
        三人都诧异问道：“这 是为何？” 宝钗道：“断断使不得。}
}
{
    \eA{Ever since that day is it that I've become alive to the fact that even a
        broken lotus leaf, and a blade of withered grass are alike worth money."
        "This is, in very truth, the way wealthy and well-to-do people talk!"
        Pao-ch'ai laughed. "But notwithstanding your honourable position, young
        ladies, you really understand nothing about these concerns.}
}
{
}

\Verse{18}
{
    \zA{你们这里多少得用的人，一个个闲着没事办， 这会子我又弄个人来，叫那起人连我也看小了。 我倒替你们想出一个人来：怡红院 有个老叶妈，他就是焙茗的娘。
        那是个诚实老人家，他又合我们莺儿妈极好。 不如 把这事交与叶妈，他有不知的，不必咱们说给他，就找莺儿的娘去商量了。 那怕叶
        妈全不管，竟交与那一个，这是他们私情儿，有人说闲话也就怨不到咱们身上。 如 此一行，你们办的又公道，于事又妥当。” 李纨平儿都道：“很是。”}
}
{
    \eA{Yet, haven't you, with all your book-lore, seen anything of the passage in
        the writing of Chu Fu-tzu: 'Throw not thyself away?'" "I've read it, it's
        true," T'an Ch'un smiled, "but its object is simply to urge people to exert
        themselves; it's as much empty talk as any random arguments, and how could
        it be bodily treated as gospel?" "Chu-tzu's work all as much empty talk as
        any random arguments?" Pao-ch'ai exclaimed.}
}
{
}

\Verse{19}
{
    \zA{探春笑道： “虽如此，只怕他们见利忘义呢。” 平儿笑道：“不相干。 前日莺儿还认了叶妈做 干娘，请吃饭吃酒，两家和厚的很呢。” 探春听了，方罢了。
        又共斟酌出几个人来， 俱是他四人素昔冷眼取中的，用笔圈出。 一时婆子们来回：“大夫已去。” 将药方送上去，三人看了。 一面遣人送出外
        边去取药，监派调服，一面探春与李纨明示诸人：某人管某处，“按四季，除家中 定例用多少外，馀者任凭你们采取去取利，年终算账。”}
}
{
    \eA{"Why every sentence in it is founded on fact. You've only had the management
        of affairs in your hands for a couple of days, and already greed and
        ambition have so beclouded your mind that you've come to look upon Chu-tzu
        as full of fraud and falsehood.}
}
{
}

\Verse{20}
{
    \zA{探春笑道：“我又想起一 件事：若年终算账，归钱时自然归到账房，仍是上头又添一层管主，还在他们手心 里又剥一层皮。
        这如今我们兴出这件事，派了你们，已是跨过他们的头去了，心里 有气只说不出来，你们年终去归账，他还不捉弄你们等什么？ 再者这一年间管什么
        的，主子有一全分，他们就得半分，这是每常的旧规，人所共知的。 如今这园子是 我的新创，竟别入他们的手，每年归账，竟归到里头来才好。”}
}
{
    \eA{But when you by and bye go out into the world and see all those mighty
        concerns reeking with greed and corruption, you'll even go so far as to
        treat Confucius himself as a fraud!" "Haven't you with all your culture read
        a book like that of Chi-tzu's?" Pan Ch'un laughed.}
}
{
}

\Verse{21}
{
    \zA{宝钗笑道：“依我 说，里头也不用归账，这个多了，那个少了，倒多了事。 不如问他们谁领这一分的， 他就揽一宗事去。 不过是园里的人动用。
        我替你们算出来了，有限的几宗事，不过 是头油、胭粉、香、纸，每一位姑娘，几个丫头，都是有定例的； 再者各处苕帚、 簸箕、掸子，并大小禽鸟鹿兔吃的粮食。
        不过这几样。 都是他们包了去，不用账房 去领钱。 你算算，就省下多少来？” 平儿笑道：“这几宗虽小，一年通共算了，也 省的下四百多银子。”}
}
{
    \eA{"Chi-tzu said in bygone days 'that when one descends into the arena where
        gain and emoluments are to be got, and enters the world of planning and
        plotting, one makes light of the injunctions of Yao and Shun, and disregards
        the principles inculcated by Confucius and Mencius.'" "What about the next
        line?" Pao-ch'ai insinuated with a significant smile.}
}
{
}

\Verse{22}
{
    \zA{宝钗笑道：“却又来。 一年四百，二年八百两，打租的房子 也能多买几间，薄沙地也可以添几亩了。 虽然还有敷馀，但他们既辛苦了一年，也
        要叫他们剩些，粘补自家。 虽是兴利节用为纲，然也不可太过，要再省上二三百银 子，失了大体统，也不像。 所以这么一行，外头帐房里一年少出四五百银子，也不
        觉的很艰啬了； 他们里头却也得些小补； 这些没营生的妈妈们，也宽裕了； 园子里 花木，也可以每年滋长繁盛； 就是你们，也得了可使之物：这庶几不失大体。}
}
{
    \eA{"I now cut the text short," T'an Ch'un smilingly rejoined, "in order to
        adapt the sense to what I want to say. Would I recite the following
        sentence, and heap abuse upon my own self; is it likely I would; eh?"
        "There's nothing under the heavens that can't be turned to some use,"
        Pao-ch'ai added. "And since everything can be utilised, everything must be
        worth money.}
}
{
}

\Verse{23}
{
    \zA{若一 味要省时，那里搜寻不出几个钱来？ 凡有些馀利的，一概入了官中，那时里外怨声 载道，岂不失了你们这样人家的大体？
        如今这园里几十个老妈妈们，若只给了这个， 那剩的也必抱怨不公； 我才说的他们只供给这个几样，也未免太宽裕了。 一年竟除
        这个之外，他每人不论有馀无馀，只叫他拿出若干吊钱来，大家凑齐，单散与这些 园中的妈妈们。 他们虽不料理这些，却日夜也都在园中照料；}
}
{
    \eA{But can it be that a person gifted with such intelligence as yours can have
        had no experience in such great matters and legitimate concerns as these?"
        "You send for a person," Li Wan laughingly interposed, 'and you don't speak
        about what's right and proper, but you start an argument on learning."
        "Learning is right and proper," Pao-ch'ai answered.}
}
{
}

\Verse{24}
{
    \zA{当差之人，关门闭户， 起早睡晚，大雨大雪，姑娘们出入，抬轿子、撑船、拉冰床一应粗重活计，都是他
        们的差使：一年在园里辛苦到头，这园内既有出息，也是分内该沾带些的。 还有一 句至小的话，越发说破了：你们只顾了自己宽裕，不分与他们些，他们虽不敢明怨，
        心里却都不服，只用假公济私的，多摘你们几个果子，多掐几枝花儿，你们有冤还 没处诉呢。 他们也沾带些利息，你们有照顾不到的，他们就替你们照顾了。”}
}
{
    \eA{"If we made no allusion to learning, we'd all soon enough drift among the
        rustic herd!" The trio bandied words for a while, after which they turned
        their attention again to pertinent affairs. T'an Ch'un took up once more the
        thread of the conversation.}
}
{
}

\Verse{25}
{
    \zA{众婆子听了这个议论，又去了帐房受辖制，又不与凤姐儿去算帐，一年不过多 拿出若干吊钱来，各各欢喜异常，都齐声说：“愿意!强如出去被他们揉搓着，还
        得拿出钱来呢。” 那不得管地的，听了每年终无故得钱，更都喜欢起来，口内说： “他们辛苦收拾，是该剩些钱粘补的； 我们怎么好‘稳吃三注’呢？” 宝钗笑道：
        “妈妈们也别推辞了，这原是分内应当的。 你们只要日夜辛苦些，别躲懒纵放人吃 酒赌钱就是了。 不然，我也不该管这事。}
}
{
    \eA{"This garden of ours," she argued, "is only half as big as theirs, so if you
        double the income they derive, you will see that we ought to reap a net
        profit of four hundred taels a year. But were we also now to secure a
        contract for our surplus products, the money, we'd earn, would, of course,
        be a mere trifle and not one that a family like ours should hanker after.}
}
{
}

\Verse{26}
{
    \zA{你们也知道，我姨娘亲口嘱托我三五回， 说大奶奶如今又不得闲，别的姑娘又小，托我照看照看。 我若不依，分明是叫姨娘 操心。
        我们太太又多病，家务也忙，我原是个闲人，就是街坊邻舍，也要帮个忙儿， 何况是姨娘托我？ 讲不起众人嫌我。 倘或我只顾沽名钓誉的，那时酒醉赌输，再生
        出事来，我怎么见姨娘？ 你们那时后悔也迟了，就连你们素昔的老脸也都丢了。}
}
{
    \eA{And were we to depute two special persons (to attend to the garden), the
        least permission given by them to any one to turn anything to improper uses,
        would, since there be so many things of intrinsic value, be tantamount to a
        reckless destruction of the gifts of heaven.}
}
{
}

\Verse{27}
{
    \zA{这 些姑娘们，这么一所大花园子，都是你们照管着，皆因看的你们是三四代的老妈妈， 最是循规蹈矩，原该大家齐心顾些体统。 你们反纵放别人，任意吃酒赌博。
        姨娘听 见了，教训一场犹可，倘若被那几个管家娘子听见了，他们也不用回姨娘，竟教导 你们一场，你们这年老的反受了小的教训。}
}
{
    \eA{So would it not be preferable to select several quiet, steady and
        experienced old matrons, out of those stationed in the grounds, and appoint
        them to put them in order and look after things? Neither will there be any
        need then to make them pay any rent, or give any taxes in kind. All we can
        ask them is to supply the household with whatever they can afford during the
        year.}
}
{
}

\Verse{28}
{
    \zA{虽是他们是管家管的着你们，何如自己 存些体面，他们如何得来作践呢!所以我如今替你们想出这个额外的进益来，也为
        的是大家齐心，把这园里周全得谨谨慎慎的，使那些有权执事的看见这般严肃谨 慎，且不用他们操心，他们心里岂不敬服？ 也不枉替你们筹画些进益了。 你们去细
        细想想这话。” 众人都欢喜说：“姑娘说的很是。 从此姑娘奶奶只管放心。 姑娘奶 奶这么疼顾我们，我们再要不体上情，天地也不容了。”}
}
{
    \eA{In the first place, the garden will, with special persons to look after the
        plants and trees, naturally so improve from year to year that there won't be
        any bustle or confusion, whenever the time draws nigh to utilise the
        grounds. Secondly, people won't venture to injure or uselessly waste
        anything.}
}
{
}

\Verse{29}
{
    \zA{刚说着，只见林之孝家的进来，说：“江南甄府里家眷昨日到京，今日进宫朝 贺，此刻先遣人来送礼请安。” 说着便将礼单送上去。 探春接了，看道是：“上用
        的妆缎蟒缎十二匹。 上用杂色缎十二匹。 上用各色纱十二匹。 上用宫绸十二匹。 宫 用各色缎纱绸绫二十四匹。” 李纨探春看过，说：“用上等封儿赏他。”
        因又命人 去回了贾母。 贾母命人叫李纨、探春、宝钗等都过来，将礼物看了。 李纨收过一边， 吩咐内库上人说：“等太太回来看了再收。”}
}
{
    \eA{In the third place, the old matrons themselves will, by availing themselves
        of these small perquisites, not labour in the gardens year after year and
        day after day all for no good. Fourthly, it will in like manner be possible
        to effect a saving in the expenditure for gardeners, rockery-layers,
        sweepers and other necessary servants. And this excess can be utilised for
        making up other deficiencies.}
}
{
}

\Verse{30}
{
    \zA{贾母因说：“这甄家又不与别家相同。 上等封儿赏男人。 只怕转眼又打发女人来请安，预备下尺头。” 一语未了，果然人回：“甄府四个女人来请安。”
        贾母听了，忙命人带进来。 那四个人都是四十往上年纪，穿带之物皆比主子不大差别。 请安问好毕，贾母便命 拿了四个脚踏来。
        他四人谢了坐，等着宝钗等坐了，方都坐下。 贾母便问：“多早 晚进京的？” 四人忙起身回说：“昨儿进的京，今儿太太带了姑娘进宫请安去了，
        所以叫女人们来请安，问候姑娘们。”}
}
{
    \eA{I don't see any reason why this shouldn't be practicable!" Pao-ch'ai was
        standing below contemplating the pictures with characters suspended on the
        walls. Upon hearing these suggestions, she readily nodded her head
        assentingly and smiled. "Excellent!" she cried. "'Within three years, there
        will be no more famines and dearths.'" "What a first-rate plan!" Li Wan
        chimed in.}
}
{
}

\Verse{31}
{
    \zA{贾母笑问道：“这些年没进京，也不想到就 来。” 四人也都笑回道：“正是。 今年是奉旨唤进京的。” 贾母问道：“家眷都来 了？”
        四人回说：“老太太和哥儿、两位小姐，并别位太太，都没来； 就只太太带 了三姑娘来了。” 贾母道：“有人家没有？” 四人道：“还没有呢。” 贾母笑道：
        “你们大姑娘和二姑娘，这两家，都和我们家甚好。” 四人笑道：“正是。 每年姑 娘们有信回来说，全亏府上照看。” 贾母笑道：“什么‘照看’？}
}
{
    \eA{"This, if actually adopted, will delight the heart of Madame Wang. Pecuniary
        economies are of themselves a paltry matter; but there will be then in the
        garden those to sweep the grounds, and those whose special charge will be to
        look after them.}
}
{
}

\Verse{32}
{
    \zA{原是世交，又是 老亲，原应当的。 你们二姑娘更好，不自尊大，所以我们才走的亲密。” 四人笑道： “这是老太太过谦了。”
        贾母又问：“你这哥儿也跟着你们老太太？” 四人回说： “也跟着老太太呢。” 贾母道：“几岁了？” 又问：“上学不曾？” 四人笑说：“今 年十三岁。
        因长的齐整，老太太很疼，自幼淘气异常，天天逃学，老爷太太也不便 十分管教。” 贾母笑道：“也不成了我们家的了？ 你这哥儿叫什么名字？”}
}
{
    \eA{Besides, were the persons selected allowed to turn up an honest cash by
        selling part of the products, they will be so impelled by a sense of their
        responsibilities, and prompted by a desire of gain that there won't any
        longer be any who won't acquit themselves of their duties to the fullest
        measure." "It remained for you, miss, to put these suggestions in words,"
        P'ing Erh remarked.}
}
{
}

\Verse{33}
{
    \zA{四人道： “因老太太当作宝贝一样，他又生的白，老太太便叫作‘宝玉’。” 贾母笑向李纨 道：“偏也叫个‘宝玉’！”
        李纨等忙欠身笑道：“从古至今，同时隔代，重名的 很多。” 四人也笑道：“起了这小名儿之后，我们上下都疑惑，不知那位亲友家也 倒像曾有一个的。
        只是这十来年没进京来，却记不真了。” 贾母笑道：“那就是我 的孙子。 ――人来。”}
}
{
    \eA{"Our mistress may have entertained the idea, but it is by no means certain
        that she thought it nice on her part to give utterance to it.}
}
{
}

\Verse{34}
{
    \zA{众媳妇丫头答应了一声，走近几步，贾母笑道：“园里把咱 们的宝玉叫了来，给这四个管家娘子瞧瞧，比他们的宝玉如何。”
        众媳妇听了，忙去了，半刻，围了宝玉进来。 四人一见，忙起身笑道：“唬了 我们一跳!要是我们不进府来，倘若别处遇见，还只当我们的宝玉后赶着也进了京 呢。”
        一面说，一面都上来拉他的手，问长问短。 宝玉也笑问个好。 贾母笑道：“比 你们的长的如何？” 李纨等笑道：“四位妈妈才一说，可知是模样儿相仿了。”}
}
{
    \eA{For as you, young ladies, live at present in the garden, she could not
        possibly, unable as she is to supply such additional ornaments as will make
        it more showy, contrariwise depute people to exercise authority in it, and
        to keep it in order, with a view of effecting a reduction in expenses. Such
        a proposal could never have dropped from her lips." Pao-ch'ai advanced up to
        her with alacrity.}
}
{
}

\Verse{35}
{
    \zA{贾 母笑道：“那有这样巧事。 大家子孩子们，再养的娇嫩，除了脸上有残疾十分丑的， 大概看去都是一样齐整，这也没有什么怪处。”
        四人笑道：“如今看来，模样是一 样!据老太太说，淘气也一样，我们看来，这位哥儿性情却比我们的好些。” 贾母 忙笑问怎么。
        四人笑道：“方才我们拉哥儿的手说话，便知道了。 若是我们那一位， 只说我们糊涂。 慢说拉手，他的东西我们略动一动也不依。 所使唤的人都是女孩子 们。”}
}
{
    \eA{Rubbing her face: "Open that mouth of yours wide," she laughed, "and let me
        see of what stuff your teeth and tongue are made! Ever since you put your
        foot out of bed this morning you've jabbered away up to this very moment!
        And your song has all been in one strain.}
}
{
}

\Verse{36}
{
    \zA{四人未说完，李纨姊妹等禁不住都失声笑出来。 贾母也笑道：“我们这会子 也打发人去见了你们宝玉，若拉他的手，他也自然勉强忍耐着。 不知你我这样人家
        的孩子，凭他们有什么刁钻古怪的毛病，见了外人，必是要还出正经礼数来的。 若 他不还正经礼数，也断不容他刁钻去了。 就是大人溺爱的，也因为他一则生的得人
        意儿； 二则见人礼数，竟比大人行出来的还周到，使人见了可爱可怜，背地里所以 才纵他一点子。}
}
{
    \eA{For neither have you been very complimentary to Miss Tertia, nor have you
        admitted that your mistress is, as far as wits go, so much below the mark as
        to be unable to effect suitable provision. Yet whenever Miss Tertia advanced
        any arguments, you've at once made use of endless words to join issue with
        her. This is because the plan devised by Miss Tertia was also hit upon by
        your lady Feng.}
}
{
}

\Verse{37}
{
    \zA{若一味他只管没里没外，不给大人争光，凭他生的怎样，也是该打 死的。” 四人听了，都笑道：“老太太这话正是。 虽然我们宝玉淘气古怪，有时见
        了客，规矩礼数，比大人还有趣，所以无人见了不爱，只说：‘为什么还打他？’ 殊不知他在家里无法无天，大人想不到的话偏会说，想不到的事偏会行，所以老爷
        太太恨的无法。 就是任性，也是小孩子的常情； 胡乱花费，也是公子哥儿的常情； 怕上学，也是小孩子的常情：都还治的过来。}
}
{
    \eA{But there must surely have been a reason why she couldn't carry it into
        execution. Again, as the young ladies have now their quarters in the garden,
        she couldn't, with any decency, direct any one to go and rule over it, for
        the mere sake of saving a few cash. Just consider this.}
}
{
}

\Verse{38}
{
    \zA{第一，天生下来这一种刁钻古怪的脾 气，如何使得？” 一语未了，人回：“太太回来了。” 王夫人进来，问过安，他四
        人请了安，大概说了两句，贾母便命：“歇歇去罢。” 王夫人亲捧过茶，方退出去。 四人告辞了贾母，便往王夫人处来，说了一会子家务，打发他们回去，不必细说。
        这里贾母喜得逢人便告诉：也有一个宝玉，也都一般行景。 众人都想着天下的 世宦人家，同名的这也很多，祖母溺爱孙子也是常事，不是什么罕事，皆不介意。}
}
{
    \eA{If the garden is actually handed to people to make profit out of it, the
        parties interested will, of course, not even permit a single spray of
        flowers to be plucked, and not a single fruit to be taken away. With such as
        come within the category of senior young ladies, they won't naturally have
        the audacity to be particular; but they'll daily have endless rows with the
        junior girls.}
}
{
}

\Verse{39}
{
    \zA{独宝玉是个迂阔呆公子的心性，自为是那四人承悦贾母之词。 后至园中去看湘云病 去，湘云因说他：“你放心闹罢，先还‘单丝不成线，独树不成林’，如今有了个
        对子了。 闹利害了，再打急了，你好逃到南京找那个去。” 宝玉道：“那里的谎话， 你也信了？ 偏又有个宝玉了？”
        湘云道：“怎么列国有个蔺相如，汉朝又有个司马 相如呢？” 宝玉笑道：“这也罢了，偏又模样儿也一样，这也是有的事吗？” 湘云
        道：“怎么匡人看见孔子，只当是阳货呢？”}
}
{
    \eA{(Lady Feng) has, with her fears about the future and her misgivings about
        the present, shown herself neither too overbearing nor too servile. This
        mistress of theirs is not friendly disposed towards us, but when she hears
        of her various proposals, shame might induce her to turn over a new leaf."}
}
{
}

\Verse{40}
{
    \zA{宝玉笑道：“孔子阳货虽同貌，却不 同名； 蔺与司马虽同名，而又不同貌。 偏我和他就两样俱同不成？” 湘云没了话答
        对，因笑道：“你只会胡搅，我也不和你分证。 有也罢，没也罢，与我无干！” 说 着，便睡下了。 宝玉心中便又疑惑起来：若说必无，也似必有；
        若说必有，又并无目睹。 心中 闷闷，回至房中榻上，默默盘算，不觉昏昏睡去，竟到一座花园之内。 宝玉诧异道： “除了我们大观园，竟又有这一个园子？”}
}
{
    \eA{"Early this morning," T'an Ch'un laughingly observed, "I was very cross, but
        as soon as I heard of her (P'ing Erh's) arrival, I casually remembered that
        her mistress employed, during her time, such domestics as were up to all
        kinds of larks, and at the sight of her, I got more cross than ever.}
}
{
}

\Verse{41}
{
    \zA{正疑惑间，忽然那边来了几个女孩儿， 都是丫鬟，宝玉又诧异道：“除了鸳鸯、袭人、平儿之外，也竟还有这一干人？” 只见那些丫鬟笑道：“宝玉怎么跑到这里来？”
        宝玉只当是说他，忙来陪笑说道： “因我偶步到此，不知是那位世交的花园？ 姐姐们带我逛逛。” 众丫鬟都笑道：“原 来不是咱们家的宝玉。
        他生的也还干净，嘴儿也倒乖觉。” 宝玉听了，忙道：“姐 姐们这里，也竟还有个宝玉？”}
}
{
    \eA{But, little though one would have thought it, she behaved from the moment
        she came, like a rat that tries to get out of the way of a cat. And as she
        had had to stand for ever so long, I pitied her very much; but she took up
        the thread of the conversation, and went on to spin that long yarn of hers.}
}
{
}

\Verse{42}
{
    \zA{丫鬟们忙道：“‘宝玉’二字，我们家是奉老太太、 太太之命，为保佑他延年消灾，我们叫他，他听见喜欢； 你是那里远方来的小厮，
        也乱叫起来!仔细你的臭肉，不打烂了你的。” 又一个丫鬟笑道：“咱们快走罢， 别叫宝玉看见。” 又说：“同这臭小子说了话，把咱们熏臭了。” 说着一径去了。
        宝玉纳闷道：“从来没有人如此荼毒我，他们如何竟这样的？ 莫不真也有我这样一 个人不成？” 一面想，一面顺步早到了一所院内。}
}
{
    \eA{Yet, instead of mentioning that her mistress treats me with every
        consideration, she, on the contrary, observed: 'The kindness with which you
        have all along dealt with our lady miss, has not been to no purpose.' This
        remark therefore not only dispelled my anger, but filled me with so much
        shame that I began to feel sore at heart.}
}
{
}

\Verse{43}
{
    \zA{宝玉诧异道：“除了怡红院，也竟还有这 么一个院落？” 忽上了台阶，进入屋内，只见榻上有一个人卧着，那边有几个女儿 做针线，或有嬉笑玩耍的。
        只见榻上那个少年叹了一声，一个丫鬟笑问道：“宝玉， 你不睡，又叹什么？ 想必为你妹妹病了，你又胡愁乱恨呢。” 宝玉听说，心下也便
        吃惊，只见榻上少年说道：“我听见老太太说，长安都中也有个宝玉，和我一样的 性情，我只不信。}
}
{
    \eA{And, when I came to think carefully over the matter, I failed to see how I,
        a mere girl, who had personally done so much mischief that not a soul cared
        a straw for me and not a soul took any interest in me, could possess any
        such good qualities as to treat any one kindly…." When she reached this
        point, she could not check her tears from brimming over. Li Wan and her
        associates perceived how pathetically she spoke;}
}
{
}

\Verse{44}
{
    \zA{我才做了一个梦，竟梦中到了都中一个大花园子里头，遇见几个 姐姐，都叫我臭小厮，不理我。 好容易找到他房里，偏他睡觉，空有皮囊，真性不 知往那里去了。”
        宝玉听说，忙说道：“我因找宝玉来到这里，原来你就是宝玉？” 榻上的忙下来拉住，笑道：“原来你就是宝玉!这可不是梦里了？” 宝玉道：“这 如何是梦？
        真而又真的！” 一语未了，只见人来说：“老爷叫宝玉。” 吓得二人皆 慌了，一个宝玉就走。 一个便忙叫：“宝玉快回来!宝玉快回来！”}
}
{
    \eA{and, recalling to mind how Mrs. Chao had always run her down, and how she
        had ever been involved in some mess or other with Madame Wang, on account of
        this Mrs. Chao, they too found it difficult to refrain from melting into
        sobs. But they then used their joint efforts to console her.}
}
{
}

\Verse{45}
{
    \zA{袭人在旁听他梦中自唤，忙推醒他，笑问道：“宝玉在那里？” 此时宝玉虽醒， 神意尚自恍惚，因向门外指说：“才去不远。” 袭人笑道：“那是你梦迷了。 你揉
        眼细瞧，是镜子里照的你的影儿。” 宝玉向前瞧了一瞧，原是那嵌的大镜对面相照， 自己也笑了。 早有丫鬟捧过漱盂茶卤来漱了口。
        麝月道：“怪道老太太常嘱咐说： ‘小人儿屋里不可多有镜子，人小魂不全，有镜子照多了，睡觉惊恐做胡梦。’}
}
{
    \eA{"Let's avail ourselves of this quiet day," they suggested, "to try and find
        out how we could increase our revenue and remove abuses, so as not to render
        futile the charge laid on us by Madame Wang. What use or purpose is it to
        allude to such trivial matters?" "I've already grasped your object," P'ing
        Erh hastily ventured. "Miss, speak out; who do you consider fit?}
}
{
}

\Verse{46}
{
    \zA{如 今倒在大镜子那里安了一张床!有时放下镜套还好，往前去天热困倦，那里想的到 放他？ 比如方才就忘了，自然先躺下照着影儿玩来着，一时合上眼自然是胡梦颠倒
        的。 不然，如何叫起自己的名字来呢？ 不如明日挪进床来是正经。” 一语未了，只 见王夫人遣人来叫宝玉。 不知有何话说，且听下回分解。 (本章完)}
}
{
    \eA{And as soon as the proper persons have been fixed upon, everything will be
        square enough." "What you say is all very well," T'an Ch'un rejoined, "but
        it will be necessary to let your lady know something about it. It has never
        been the proper thing for us in here to scrape together any small profits.
        But as your mistress is full of gumption, I adopted the course I did.}
}
{
}

\Verse{47}
{
    \zA{}
}
{
    \eA{Had she been at all narrowminded, with many prejudices and many jealousies,
        I wouldn't have shown the least willingness in the matter. But, as it will
        look as if I were bent upon pulling her to pieces, how can I take action
        without consulting her?" "In that case," P'ing Erh smiled, "I'll go and tell
        her something about it." With this response, she went on the errand; and
        only returned after a long lapse of time. "I said," she laughed, "that it
        would be perfectly useless for me to go. How ever could our lady not readily
        accede to an excellent proposal like this?" Hearing this, T'an Ch'un
        forthwith joined Li Wan in directing a servant to ask for the roll,
        containing the names of the matrons in the garden, and bring it to them.}
}
{
}

\Verse{48}
{
    \zA{}
}
{
    \eA{When produced, they all held council together, and fixing cursorily upon
        several persons, they summoned them to appear before them. Li Wan then
        explained to them the general outline of their duties; and not one was there
        among the whole company, who listened to her, who would not undertake the
        charge. One said: "If you confide that bamboo tree for twelve months to my
        care, it will again next year be a single tree, but besides the shoots,
        which will have been eaten at home, I shall be able, in the course of the
        year, to also pay in some money."}
}
{
}

\Verse{49}
{
    \zA{}
}
{
    \eA{"Hand me over," another one remarked, "that portion of paddy field, and
        there will, during the year, be no need to touch any public funds on account
        of the various birds, large and small, which are kept for mere fun. Besides
        that, I shall be in a position to give in something more." T'an Ch'un was
        about to pass a remark when a servant reported that the doctor had come; and
        that he had entered the garden to see Miss Shih. So the matrons were obliged
        to go and usher the doctor in. "Were there a hundred of you here," promptly
        expostulated P'ing Erh, "you wouldn't know what propriety means! Are there
        perchance no couple of housekeepers about to push themselves forward and see
        the doctor in?" "There's dame Wu and dame T'an," the servant, who brought
        the message, replied.}
}
{
}

\Verse{50}
{
    \zA{}
}
{
    \eA{"The two are on duty at the south-west corner at the 'accumulated splendour'
        gate." At this answer, P'ing Erh allowed the subject to drop. After the
        departure of the matrons, T'an Ch'un inquired of Pao-ch'ai what she thought
        of them. "Such as are diligent at the outset," Pao-ch'ai answered smiling,
        "become remiss in the end; and those who have a glib tongue have an eye to
        gain." T'an Ch'un listened to her reply; and nodding her head, she extolled
        its wisdom. Then showing them with her finger several names on the list, she
        submitted them for the perusal of the trio. P'ing Erh speedily went and
        fetched a pen and inkslab. "This old mother Chu," the trio observed, "is a
        trustworthy woman. What's more, this old dame and her sons have generation
        after generation done the sweeping of the bamboo groves.}
}
{
}

\Verse{51}
{
    \zA{}
}
{
    \eA{So let's now place the various bamboo trees under her control. This old
        mother T'ien was originally a farmer, and everything in the way of
        vegetables and rice, in and about the Tao Hsiang village, should, albeit
        they couldn't, planted as they are as a mere pastime, be treated in such
        earnest as to call for large works and extensive plantations, be entrusted
        to her care; for won't they fare better if she can be on the spot and tend
        them with extra diligence at the proper times and seasons?" "What a pity it
        is," T'an Ch'un proceeded smilingly, "that two places so spacious as the
        Heng Wu garden and the I Hung court bring no grit to the mill." "Things in
        the Heng Wu garden are in a worse state," Li Wan hastily interposed.}
}
{
}

\Verse{52}
{
    \zA{}
}
{
    \eA{"Aren't the scented wares and scented herbs sold at present everywhere in
        perfumery shops, large fairs and great temples the very counterpart of these
        things here? So if you reckon up, you will find how much greater a return
        these articles will give than any other kind of product. As for the I Hung
        court, we needn't mention other things, but only take into account the roses
        that bud during the two seasons of spring and summer; to how many don't they
        amount in all? Besides these, we've got along the whole hedge, cinnamon
        roses and monthly roses, stock roses, honey-suckle and westeria.}
}
{
}

\Verse{53}
{
    \zA{}
}
{
    \eA{Were these various flowers dried and sold to the tea and medicine shops,
        they'd also fetch a good deal of money." "Quite so!" T'an Ch'un acquiesced
        with a smile. "The thing is that there's no one with any notion how to deal
        with scented herbs." "There's Ying Erh who waits on Miss Pao-ch'ai," P'ing
        Erh promptly smiled. "Her mother is well-versed in these things. It was only
        the other day that she plucked a few, and plaited them, after drying them
        well in the sun, into a flower-basket and a gourd, and gave them to me to
        play with. But miss can you have forgotten all about it?" "I was this very
        minute speaking in your praise," Pao-ch'ai observed smiling, "and do you
        come to chaff me?" "What makes you say so?" exclaimed the trio, in utter
        astonishment.}
}
{
}

\Verse{54}
{
    \zA{}
}
{
    \eA{"It will on no account do," Pao-ch'ai added. "You employ such a lot of
        people in here that they all lead a lazy life and have nothing to put a hand
        to, and were I also now to introduce some more, that tribe will look even
        upon me with utter contempt. But let me think of some one for you. There's
        in the I Hung court, an old dame Yeh; she's Pei Ming's mother. That woman is
        an honest old lady; and is furthermore on the best of terms with our Ying
        Erh's mother. So wouldn't it be well were this charge given to this dame
        Yeh? Should there even be anything that she doesn't know, there'll be no
        necessity for us to tell her. She can go straightway and consult with Ying
        Erh's mother. And if she can't attend to everything herself, it won't matter
        to whom she relegates some of her duties.}
}
{
}

\Verse{55}
{
    \zA{}
}
{
    \eA{These will be purely private favours. In the event too of any one making any
        mean insinuations, the blame won't fall on our shoulders. By adopting this
        course, you'll be managing things in such a way as to do extreme justice to
        all; and the trust itself will also be placed on a most satisfactory
        footing." "Excellent!" ejaculated Li Wan and P'ing Erh simultaneously. "This
        may be well and good," T'an Ch'un laughed, "but the fear is that at the
        sight of gain, they'll forget all about propriety." "That's nothing to do
        with us!" P'ing Erh rejoined a smile playing, about her lips. "It was only
        the other day that Ying Erh recognised dame Yeh as her adopted mother, and
        invited her to eat and drink with them, so that the two families are on the
        most intimate terms."}
}
{
}

\Verse{56}
{
    \zA{}
}
{
    \eA{At this assurance, T'an Ch'un relinquished the topic of conversation, and,
        holding council together, they selected several persons, all of whom the
        four had ever viewed with impartial favour and they marked off their names,
        by dotting them with a pen. In a little while, the matrons came to report
        that 'the doctor had gone;' and they handed the prescription. Their three
        mistresses then perused its contents. On the one hand, they despatched
        domestics to take it outside, so that the drugs should be got, and to
        superintend their decoction. On the other, T'an Ch'un and Li Wan explicitly
        explained to the various servants chosen what particular place each had to
        look after.}
}
{
}

\Verse{57}
{
    \zA{}
}
{
    \eA{"Exclusive," they added, "of what fixed custom requires for home consumption
        during the four seasons, you are still at liberty to pluck whatever remains
        and have it taken away. As for the profits, we'll settle accounts at the
        close of the year." "I've also bethought myself of something," T'an Ch'un
        smiled. "If the settlement of accounts takes place at the end of the year,
        the money will, at the time of delivery, be naturally paid into the
        accountancy. Those high up will then as usual add a whole lot of
        controllers; and these will, on their part, fleece their own share as soon
        as the money gets into the palms of their hand.}
}
{
}

\Verse{58}
{
    \zA{}
}
{
    \eA{But as by this system, we've now initiated, you've been singled out for
        appointment, you've already ridden so far above their heads, that they
        foster all sorts of animosity against you. They don't, however, give vent to
        their feelings; but if they don't seize the close of the year, when you have
        to deliver your accounts, to play their tricks on you, for what other
        chances will they wait? Moreover, they obtain, in everything that comes
        under their control during the year, half of every share their masters get.
        This is an old custom. Every one is aware of its existence. But this is a
        new regime I now introduce in this garden, so don't let the money find its
        way into their hands! Whenever the annual settling of accounts arrives,
        bring them in to us."}
}
{
}

\Verse{59}
{
    \zA{}
}
{
    \eA{"My idea is," Pao-ch'ai smilingly suggested, "that no accounts need be
        handed even inside. This one will have a surplus, that one a deficit, so
        that it will involve no end of trouble; wouldn't it be better therefore if
        we were to find out who of them would take over this or that particular kind
        and let them purvey the various things? These are for the exclusive use of
        the inmates of the garden; and I've already made an estimate of them for
        you. They amount to just a few sorts, and simply consist of head-oil, rouge,
        powder and scented paper; in all of which, the young ladies and maids are
        subject to a fixed rule. Then, besides these, there are the brooms,
        dust-baskets and poles, wanted in different localities, and the food for the
        large and small animals and birds, and the deer and rabbits.}
}
{
}

\Verse{60}
{
    \zA{}
}
{
    \eA{These are the only kinds of things required. And if they contract for them,
        there'll be little need for any one to go to the accountancy for money. But
        just calculate what a saving will thus be effected!" "All these items are, I
        admit, mere trifles," P'ing Erh smiled, "but if you lump together what's
        used during a year, you will find that a saving of four hundred taels will
        be effected." "Again!" smilingly remarked Pao-ch'ai, "it would be four
        hundred taels in one year; but eight hundred taels in two years; and with
        these, we could purchase a few more houses and let them; and in the way of
        poor, sandy land we could also add several acres to those we've already got.
        'There will, of course, still remain a surplus;}
}
{
}

\Verse{61}
{
    \zA{}
}
{
    \eA{but as they will have ample trouble and inconvenience to put up with during
        the year, they should also be allowed some balance in hand so as to make up
        what's wanted for themselves. The main object is, of course, to increase
        profits and curtail expenses, yet we couldn't be stingy to any excessive
        degree. In fact, were we even able to make any further economy of over two
        or three hundred taels, it would never be the proper thing; should this
        involve a breach of the main principles of decorum. With this course duly
        put into practice, outside, the accountancy will issue in one year four or
        five hundred taels less, without even the semblance of any parsimony;}
}
{
}

\Verse{62}
{
    \zA{}
}
{
    \eA{while, inside, the matrons will obtain, on the other hand, some little thing
        to supply their wants with; the nurses, who have no means of subsistence,
        will likewise be placed in easy circumstances; and the plants and trees in
        the garden will year by year increase in strength and grow more abundantly.
        In this wise, you too will have such articles as will be fit for use. So
        that this plan will, to some extent, not constitute a breach of the high
        principles of propriety. And if ever we want to retrench a little more from
        where won't we be able to get money? But if the whole balance, if any, be
        put to the credit of the public fund, every one, inside as well as outside,
        will fill the streets with the din of murmurings!}
}
{
}

\Verse{63}
{
    \zA{}
}
{
    \eA{And won't this be then a slur upon the code of honour of a household such as
        yours? So were any charge to be entrusted to this one, out of the several
        tens of old nurses at present employed in the garden, and not to that one,
        the remainder will naturally resent such injustice. As I said a while back
        all that these women will have to provide among themselves amounts to a few
        articles, so they will unavoidably have ample means.}
}
{
}

\Verse{64}
{
    \zA{}
}
{
    \eA{Hence each should be told to contribute, beyond the articles that fall to
        her share during the year, a certain number of tiaos, whether she may or may
        not realise any balance, and then jointly lump these sums together, and
        distribute them among those nurses only on service in the garden. For
        although they may not have anything to do with the control of these things,
        they themselves will have to stay in the grounds, to keep an eye over the
        servants on duty, to shut the doors, to close the windows and to get up
        early and retire late. Whenever it rains in torrents or it snows hard and
        chairs have to be carried, for you, young ladies, to go out and come in;}
}
{
}

\Verse{65}
{
    \zA{}
}
{
    \eA{or boats have to be punted, and sledges drawn, these rough and arduous
        duties come alike within their sphere of work. They have to labour in the
        garden from one year's end to the other, and though, they earn something in
        those grounds, it's only right that they should able to get some small
        benefits in the discharge of their legitimate duties. But there's another
        most trivial point that I would broach with less reserve. If you only think
        of your ease, and don't share the profits with them, they will, of course,
        never presume to show their displeasure, but in their hearts they won't
        cherish you any good feeling. What they'll do will be to make public
        business a pretext to serve their own private ends with; they'll pluck more
        of your fruits than they should;}
}
{
}

\Verse{66}
{
    \zA{}
}
{
    \eA{and cut greater quantities of your flowers than they ought. And you people
        will have a grievance, but you won't have anywhere to go and confide it. But
        should they too reap some gain, they'll readily look after such things on
        your behalf as you won't have the time to attend to." The matrons listened
        to her explanations; (and finding that) they would be removed from the
        control of the accountancy, that they would not be compelled to go and
        settle accounts with lady Feng, and that all that they would be called upon
        to do every year would be to supply a few more tiaos, were each and all
        delighted to an exceptional degree. So much so, that every one of them
        exclaimed in a chorus that they were quite prepared to agree to the terms.}
}
{
}

\Verse{67}
{
    \zA{}
}
{
    \eA{"It is better," they said, "than to be obliged to go out and be squeezed by
        them; and to have to fork out our own money as well." Those too not
        entrusted with the care of any portion of land were also highly elated, when
        they heard that at the close of each year they would, though they had no
        valid claim, come in for some share of hard cash. "They'll have to bear the
        trouble," they however argued, "to keep things in order, so it's only right
        that they should be left with a few cash to meet their various wants with;
        and how could we very well gobble our three meals without doing a stroke of
        work?" "Worthy dames," Pao-ch'ai smiled, "you mustn't decline. These duties
        are within your province and you should fulfil them.}
}
{
}

\Verse{68}
{
    \zA{}
}
{
    \eA{All you need do is to exert yourselves a bit by day and night, and not be so
        remiss and careless as to suffer any of the servants to drink and gamble;
        that's all. Otherwise, I myself must have nothing to do with the control.
        But you, yourselves, know well enough that it's my aunt who appealed to me
        with her own lips three and five times to do it as a favour to her. 'Your
        eldest sister-in-law,' she represented, 'has at present no leisure, and the
        other girls are young,' and then she asked me to look after things. So if I
        now don't accede, it's as clear as day that I shall be the cause of much
        worry to my aunt. Our lady Feng herself is seriously ill, and our domestic
        affairs can't hang fire.}
}
{
}

\Verse{69}
{
    \zA{}
}
{
    \eA{I'm really with nothing to do, so were even a mere neighbour to solicit my
        help, I would also feel bound to lend her a hand in her pressure of work.
        How much more therefore when it's my own aunt, who invokes my aid? Setting
        aside the way I'm execrated by one and all, how would I ever be able to
        stare my aunt in the face, if, while I gave my sole mind to winning fame and
        fishing for praise, any one got so intoxicated and lost so much in gambling
        as to stir up trouble? At such a juncture remorse on your part will be too
        late! Even the old reputation you have ever enjoyed will entirely be lost
        and gone.}
}
{
}

\Verse{70}
{
    \zA{}
}
{
    \eA{Those young ladies and girls and this vast garden are alike placed under
        your supervision, purely and simply because one takes into account that you
        have been nurses to three or four generations and that you have most
        scrupulously observed the rules of etiquette and propriety. It's but fair
        that you should try, with one mind, and show some little regard for what's
        right and proper. But if you contrariwise behave with such laxity as to let
        people gratify their wishes by guzzling and gambling, and my aunt comes to
        hear of these nice doings, a little scolding from her will be of little
        consequence. But if the various women, who attend to the household, get
        scent of the state of affairs, they will haul you over the coals, without
        even so much as breathing one single word beforehand to my aunt.}
}
{
}

\Verse{71}
{
    \zA{}
}
{
    \eA{And venerable people, though you are, you will then, instead of tendering
        advice to young people, be called to account by them. As housekeepers, they
        exercise, it's true, authority over you; but why shouldn't you yourselves
        observe a certain amount of decorum? And if you do so, will they have any
        occasion to bully you? The reason why I've now bethought myself of this
        special boon for you is that you should unanimously strain every nerve to
        diligently attend to the garden, in order that the powers that be may, at
        the sight of your unrelenting care and zeal, have no cause to give way to
        solicitude. And won't they inwardly look up to you with regard? Neither will
        you render of no effect the various benefits devised for them.}
}
{
}

\Verse{72}
{
    \zA{}
}
{
    \eA{But go now and minutely ponder over all my advice!" All the women received
        her words with gratification. "What you say is quite right," they replied.
        "From this time forth you, miss, and you, our lady, can well compose your
        minds. With the interest both of you feel on our behalf, may heaven and
        earth not spare us, if we do not display a full amount of gratitude for all
        your kindnesses." These assurances were still being uttered when they saw
        Lin Chih-hsiao's wife walk in. "The family of the Chen mansion of Chiang
        Nan," she explained, "arrived in the capital yesterday. To-day, they're
        going into the palace to offer their congratulations. But they've now sent
        messengers ahead to come and bring presents and pay their respects." While
        she spoke, she produced the list of presents and handed it up.}
}
{
}

\Verse{73}
{
    \zA{}
}
{
    \eA{T'an Ch'un took it over from her. "They consist," she said, perusing it, "of
        twelve rolls of brocades and satins embroidered with dragons, such as are
        for imperial use; twelve rolls of satins of various colours, of the kind
        worn by the Emperor; twelve rolls of every sort of imperial gauze; twelve
        rolls of palace silks of the quality used by his majesty; and twenty rolls
        of satins, gauzes, silks and thin silks of different colours, generally worn
        by officials." After glancing over the list, Li Wan and T'an Ch'un suggested
        that a first-class tip should be given to the messengers who brought them,
        after which, they went on to direct a servant to convey the tidings to
        dowager lady Chia. Old lady Chia gave orders to call Li Wan, T'an Ch'un,
        Pao-ch'ai and the other girls.}
}
{
}

\Verse{74}
{
    \zA{}
}
{
    \eA{On their arrival, the presents were passed under review; and this over, Li
        Wan put them aside. "You must wait," she said to the servants of the inner
        store-room, "until Madame Wang comes back and sees them; you can then lock
        them up." "This Chen family too," old lady Chia thereupon added, "isn't like
        any other family; the highest tips should therefore be conferred upon the
        men. But as in a twinkle, they may also send some of their womankind to come
        and make their obeisance, silks should be got ready in anticipation."
        Scarcely was this remark concluded before a domestic actually announced:
        'that four ladies of the Chen mansion had come to pay their respects.' Upon
        hearing this, dowager lady Chia hastily directed that they should be
        introduced into her presence.}
}
{
}

\Verse{75}
{
    \zA{}
}
{
    \eA{The four women ranged from forty years and over. Their clothing and
        head-gear were not, in any material degree, different from those of
        mistresses. As soon as they presented their compliments and inquired about
        their healths, old lady Chia desired that four footstools should be moved
        forward. But though the four women thanked her for bidding them sit down,
        they only occupied the stools, after Pao-ch'ai had seated herself. "When did
        you enter the capital?" old lady Chia inquired. The four women jumped to
        their feet with alacrity. "We entered the capital yesterday," they answered.
        "Our lady has taken our young lady today into the palace to pay their
        homage.}
}
{
}

\Verse{76}
{
    \zA{}
}
{
    \eA{That's why she bade us come and give you their compliments, and see how the
        young ladies are getting on." "You hadn't paid a visit to the capital for
        ever so many years," dowager lady Chia smilingly observed, "and here you
        appear now quite unexpectedly!" The four women simultaneously smiled again.
        "Quite so!" they said. "We received this year imperial orders, summoning us
        to the capital!" "Has the whole family come?" old lady Chia asked. "Our old
        mistress, our young master, the two young ladies and the other ladies
        haven't come up," the four women explained. "Only our lady has come,
        together with Miss Tertia." "Is she engaged to any one?" old lady Chia
        asked. "Not yet," rejoined the quartet.}
}
{
}

\Verse{77}
{
    \zA{}
}
{
    \eA{"The two families, that of your senior married lady and that of your lady
        Secunda are both on most intimate terms with ours," dowager lady Chia
        smilingly added. "Yes, they are," replied the four women with a smile. "The
        letters received each year from our young ladies, assure us that they're
        entirely dependent upon the kindness bestowed upon them, in your worthy
        mansion, for their well-being." "What kindness?" old lady Chia exclaimed
        laughingly. "These two families are really friends of long standing. In
        addition to this, they're old relatives. So what we do is our simple bounden
        duty. What's more in the favour of your two young ladies is, that they're
        not full of their own importance. That's how it is that we've come to be on
        such close terms." The four women smiled.}
}
{
}

\Verse{78}
{
    \zA{}
}
{
    \eA{"This is mainly due to your venerable ladyship's excessive humility," they
        answered. "Is that young gentleman of yours too with your old mistress?" old
        lady Chia went on to inquire. "Yes, he has also come with our old mistress,"
        the four women retorted. "How old is he?" old lady Chia then asked. "Does he
        go to school?" she afterwards inquired. "He's thirteen this year," the four
        women said by way of response. "But all through those good looks of his, our
        old mistress cherishes him so fondly that from his youth up, he has been
        wayward to the extreme, and that he now daily plays the truant. But our
        master and mistress as well don't keep any great check over him." "Yet, he
        can't resemble that young fellow of ours," old lady Chia laughed. "What's
        the name of your young gentleman?"}
}
{
}

\Verse{79}
{
    \zA{}
}
{
    \eA{"As our old mistress treats him just like a real precious gem," the quartet
        explained, "and as his complexion is naturally so white, her ladyship calls
        him Pao-yü." "Here's another one with the name of Pao-yü!" old lady Chia
        laughingly said to Li Wan. Li Wan and her companions hastily made a curtsey.
        "There have been, from old times to the present," they smiled, "very many
        among contemporaries and persons of different generations as well, who have
        borne duplicate names." The four women also smiled.}
}
{
}

\Verse{80}
{
    \zA{}
}
{
    \eA{"After the selection of this infant name," they proceeded, "we all, both
        high or low, began to give way to surmises, as we could not make out in what
        relative's or friend's family there was a lad also called by the same name.
        But as we hadn't come to the capital for ten years or so, we couldn't
        remember." "That young fellow is my grandson," dowager lady Chia remarked.
        "Hallo! some one come here!" The married women and maids assented and
        approached several steps. "Go into the garden," old lady Chia smilingly
        said, "and call our Pao-yü here, so that these four housekeeping dames
        should see how he compares with their own Pao-yü." The married women, upon
        hearing her orders, promptly went off. After a while, they entered the room
        pressing round Pao-yü.}
}
{
}

\Verse{81}
{
    \zA{}
}
{
    \eA{The moment the four dames caught sight of him, they speedily rose to their
        feet. "He has given us such a start!" they exclaimed smilingly. "Had we not
        come into your worthy mansion, and perchance, met him, elsewhere, we would
        have taken him for our own Pao-yü, and followed him as far as the capital."
        While speaking they came forward and took hold of his hands and assailed him
        with questions. Pao-yü however also put on a smile and inquired after their
        healths. "How do his looks compare with those of your young gentleman?"
        dowager lady Chia asked as she smiled. "The way the four dames ejaculated
        just now," Li Wan and her companions explained, "was sufficient to show how
        much they resemble in looks." "How could there ever he such a coincidence?"
        old lady Chia laughed.}
}
{
}

\Verse{82}
{
    \zA{}
}
{
    \eA{"Yet, the children of wealthy families are so delicately nurtured that
        unless their faces are so deformed as to make them downright ugly, they're
        all equally handsome, as far as general appearances go. So there's nothing
        strange in this!" "As we gaze at his features," the quartet added, with
        smiling faces, "we find him the very image of him; and from what we gather
        from your venerable ladyship, he's also like him in waywardness. But, as far
        as we can judge, this young gentleman's disposition is ever so much better
        than that of ours." "What makes you think so?" old lady Chia precipitately
        inquired. "We saw it as soon as we took hold of the young gentleman's
        hands," the four women laughingly rejoined, "and when he spoke to us.}
}
{
}

\Verse{83}
{
    \zA{}
}
{
    \eA{Had it been that fellow of ours, he would have simply called us fools. Not
        to speak of taking his hand in ours, why we daren't even slightly move any
        of his things. That's why, those who wait on him are invariably young
        girls." Before the four dames had time to conclude what they had to say, Li
        Wan and the rest found it so hard to check themselves that with one voice
        they burst into loud laughter. Old lady Chia also laughed. "Let's also send
        some one now," she said, "to have a look at your Pao-yü. When his hand is
        taken, he too is sure to make an effort to put up with it. But don't you
        know that children of families such as yours and mine are bound,
        notwithstanding their numerous perverse and strange defects, to return the
        orthodox civilities, when they come across any strangers.}
}
{
}

\Verse{84}
{
    \zA{}
}
{
    \eA{But should they not return the proper civilities, they should, by no manner
        of means, be suffered to behave with such perverseness. It's the way that
        grown-up people doat on them that makes them what they are. And as they can,
        first and foremost, boast of bewitching good looks and they comport
        themselves, secondly, towards visitors with all propriety—, in fact, with
        less faulty deportment than their very seniors—, they manage to win the love
        and admiration of such as only get a glimpse of them. Hence it is that
        they're secretly indulged to a certain degree.}
}
{
}

\Verse{85}
{
    \zA{}
}
{
    \eA{But if they don't show the least regard to any one inside or outside, and so
        reflect no credit upon their parents, they deserve, with all their handsome
        looks, to be flogged to death." These sentiments evoked a smile from the
        four dames. "Your words venerable lady," they exclaimed, "are quite correct.
        But though our Pao-yü be wilful and strange in his ways, yet, whenever he
        meets any visitors, he behaves with courteousness and good manners; so much
        so, that he's more pleasing to watch than even grown-up persons. There is no
        one, therefore, who sees him without falling in love with him. But you'll
        say: 'why is he then beaten?' You really aren't aware that at home he has no
        regard either for precept or for heaven;}
}
{
}

\Verse{86}
{
    \zA{}
}
{
    \eA{that he comes out with things that never suggest themselves to the
        imagination of grown-up people, and that he does everything that takes one
        by surprise. The result is that his father and mother are driven to their
        wits' ends. But wilfulness is natural to young children. Reckless
        expenditure is a common characteristic of young men. Antipathy to school is
        a common feeling with young people. Yet there are ways and means to bring
        him round. The worse with him is that his disposition is so crotchety and
        whimsical. Can this ever do?…." This reply was barely ended when a servant
        informed them that their mistress had returned. Madame Wang entered the
        room, and saluted the women. The four dames paid their obeisance to her.}
}
{
}

\Verse{87}
{
    \zA{}
}
{
    \eA{But they had just had sufficient time to pass a few general observations,
        when dowager lady Chia bade them go and rest. Madame Wang then handed the
        tea in person and withdrew from the apartment. But when the four dames got
        up to say good-bye, old lady Chia adjourned to Madame Wang's quarters. After
        a chat with her on domestic affairs, she however told the women to go back;
        so let us put them by without any further allusion to them. During this
        while, old lady Chia's spirits waxed so high, that she told every one and
        any one she came across that there was another Pao-yü, and that he was, in
        every respect, the very image of her grandson.}
}
{
}

\Verse{88}
{
    \zA{}
}
{
    \eA{But as each and all bore in mind that there were many inmates among the
        large households of those officials with official ancestors, called by the
        same names, that it was an ordinary occurrence for a grandmother to be
        passionately fond of her grandson, and that there was nothing out-of-the-way
        about it, they treated the matter as of no significance. Pao-yü alone
        however was such a hair-brained simpleton that he conjectured that the
        statements made by the four dames had been intended to flatter his
        grandmother Chia. But subsequently he betook himself into the garden to see
        how Shih Hsiang-yün was getting on. "Compose your mind now," Shih Hsiang-yün
        then said to him, "and go on with your larks!}
}
{
}

\Verse{89}
{
    \zA{}
}
{
    \eA{Once, you were as lonely as a single fibre, which can't be woven into
        thread, and like a single bamboo, which can't form a grove, but now you've
        found your pair. When you exasperate your parents, and they give you beans,
        you'll be able to bolt to Nanking in quest of the other Pao-yü." "What utter
        rubbish!" Pao-yü exclaimed. "Do you too believe that there's another
        Pao-yü?" "How is it," Hsiang-yün asked, "that there was some one in the Lieh
        state called Lin Hsiang-ju, and that during the Han dynasty there lived
        again another person, whose name was Ssu Ma Hsiang-ju?" "This matter of
        names is all well enough," Pao-yü rejoined with a smile. "But as it happens,
        his very appearance is the counterpart of mine. Such a thing could never
        be!"}
}
{
}

\Verse{90}
{
    \zA{}
}
{
    \eA{"How is it," Hsiang-yün inquired, "that when the K'uang people saw
        Confucius, they fancied it was Yang Huo?" "Confucius and Yang Huo," Pao-yü
        smilingly argued, "may have been alike in looks, but they hadn't the same
        names. Lin and Ssu were again, notwithstanding their identical names,
        nothing like each other in appearances. But can it ever be possible that he
        and I should resemble each other in both respects?" Hsiang-yün was at a loss
        what reply to make to his arguments. "You may," she consequently remarked
        smiling, "propound any rubbish you like, I'm not in the humour to enter into
        any discussion with you. Whether there be one or not is quite immaterial to
        me. It doesn't concern me at all." Saying this, she lay herself down. Pao-yü
        however began again to exercise his mind with further surmises.}
}
{
}

\Verse{91}
{
    \zA{}
}
{
    \eA{"If I say," he cogitated, "that there can't be one, there seems from all
        appearances to be one. And if I say that there is one, I haven't, on the
        other hand, seen him with my own eyes." Sad and dejected he returned
        therefore to his quarters, and reclining on his couch, he silently communed
        with his own thoughts until he unconsciously became drowsy and fell fast
        asleep. Finding himself (in his dream) in some garden or other, Pao-yü was
        seized with astonishment. "Besides our own garden of Broad Vista," he
        reflected, "is there another such garden?" But while indulging in these
        speculations, several girls, all of whom were waiting-maids, suddenly made
        their appearance from the opposite direction. Pao-yü was again filled with
        surprise.}
}
{
}

\Verse{92}
{
    \zA{}
}
{
    \eA{"Besides Yüan Yang, Hsi Jen and P'ing Erh," he pondered, "are there verily
        such maidens as these?" "Pao-yü!" he heard that company of maids observe,
        with faces beaming with smiles, "how is it you find yourself in here?"
        Pao-yü laboured under the impression that they were addressing him. With
        hasty step, he consequently drew near them, and returned their smiles. "I
        got here," he answered, "quite listlessly. What old family friend's garden
        is this, I wonder? But sisters, pray, take me for a stroll." The maids
        smiled with one consent. "Really!" they exclaimed, "this isn't our Pao-yü.
        But his looks too are spruce and nice; and he is as precocious too with his
        tongue." Pao-yü caught their remarks. "Sisters!" he eagerly cried, "is there
        actually a second Pao-yü in here?"}
}
{
}

\Verse{93}
{
    \zA{}
}
{
    \eA{"As for the two characters 'Pao-yü,'" the maids speedily explained, "every
        one in our house has received our old mistress' and our mistress'
        injunctions to use them as a spell to protract his life for many years and
        remove misfortune from his path, and when we call him by that name, he
        simply goes into ecstasies, at the very mention of it. But you, young brat,
        from what distant parts of the world do you hail that you've recklessly been
        also dubbed by the same name? But beware lest we pound that frowzy flesh of
        yours into mincemeat." "Let's be off at once!" urged another maid, as she
        smiled. "Don't let our Pao-yü see us here and say again that by hobnobbing
        with this stinking young fellow, we've been contaminated by all his
        pollution."}
}
{
}

\Verse{94}
{
    \zA{}
}
{
    \eA{With these words on her lips, they straightway walked off. Pao-yü fell into
        a brown study. "There's never been," he mused, "any one to treat me with
        such disdain before! But what is it, in fact, that induces them to behave
        towards me in this manner? May it not be true that there lives another human
        being the very image of myself?" While lost in reverie, he advanced with
        heedless step, until he reached a courtyard. Pao-yü was struck with wonder.
        "Is there actually," he cried, "besides the I Hung court another court like
        it?" Spontaneously then ascending the steps, he entered an apartment, in
        which he discerned some one reclining on a couch. On the off side sat
        several girls, busy at needlework; now laughing joyfully; now practising
        their jokes; when he overheard the young person on the couch heave a sigh.}
}
{
}

\Verse{95}
{
    \zA{}
}
{
    \eA{"Pao-yü," smilingly inquired a maid, "what, aren't you asleep? What are you
        once more sighing for? I presume it's because your sister is ill that you
        abandon yourself again to idle fears and immoderate anguish!" These words
        fell on Pao-yü's ears, and took him quite aback. "I've heard grandmother
        say," he overheard the young person on the couch observe, "that there lives
        at Ch'ang An, the capital, another Pao-yü endowed with the same disposition
        as myself. I never believed what she told me; but I just had a dream, and in
        this dream I found myself in a garden of the metropolis where I came across
        several maidens;}
}
{
}

\Verse{96}
{
    \zA{}
}
{
    \eA{all of whom called me a 'stinking young brat,' and would have nothing
        whatever to do with me. But after much difficulty, I succeeded in
        penetrating into his room. He happened to be fast asleep. There he lay like
        a mere bag of bones. His real faculties had flown somewhere or other;
        whither it was hard for me to say." Hearing this, "I've come here," Pao-yü
        said with alacrity, "in search of Pao-yü; and are you, indeed, that Pao-yü?"
        The young man on the couch jumped down with all haste and enfolded him in
        his arms. "Are you verily Pao-yü?" he laughingly asked. "This isn't by any
        means such stuff as dreams are made of!" "How can you call this a dream?"
        Pao-yü rejoined. "It's reality, yea, nothing but reality!"}
}
{
}

\Verse{97}
{
    \zA{}
}
{
    \eA{But scarcely was this rejoinder over, than he heard some one come, and say:
        "our master, your father, wishes to see you, Pao-yü." The two lads started
        with fear. One Pao-yü rushed off with all despatch. The other promptly began
        to shout, "Pao-yü! come back at once! Pao-yü; be quick and return!" Hsi Jen,
        who stood by (Pao-yü), heard him call out his own name, in his dreams, and
        immediately gave him a push and woke him up. "Where is Pao-yü gone to?" she
        laughed. Although Pao-yü was by this time aroused from sleep, his senses
        were as yet dull, so pointing towards the door, "He's just gone out," he
        replied, "he's not far off." Hsi Jen laughed. "You're under the delusion of
        a dream," she said.}
}
{
}

\Verse{98}
{
    \zA{}
}
{
    \eA{"Rub your eyes and look carefully! It's your reflection in the mirror."
        Pao-yü cast a glance in front of him, and actually caught sight of the large
        inlaid mirror, facing him quite opposite, so he himself burst out laughing.
        But, presently, a maid handed him a rince-bouche and tea and salt, and he
        washed his mouth. "Little wonder is it," She Yüeh ventured, "if our old
        mistress has repeatedly enjoined that it isn't good to have too many mirrors
        about in young people's rooms, for as the spirit of young persons is not
        fully developed there is every fear, with mirrors casting their reflections
        all over the place, of their having wild dreams in their sleep. And is a bed
        now placed before that huge mirror there?}
}
{
}

\Verse{99}
{
    \zA{}
}
{
    \eA{When the covers of the mirrors are let down, no harm can befall; but as the
        season advances, and the weather gets hot, one feels so languid and tired,
        that is one likely to think of dropping them? Just as it happened a little
        time back; it slipped entirely from your memory. Of course, when he first
        got into bed, he must have played with his face towards the glass; but upon
        shortly closing his eyes, he must naturally have fallen into such confused
        dreams, that they thoroughly upset his rest. Otherwise, how is it possible
        that he should have started shouting his own name? Would it not be as well
        if the bed were moved inside to-morrow? That's the proper place for it."
        Hardly had she, however, done, before they perceived a servant, sent by
        Madame Wang to call Pao-yü.}
}
{
}

\Verse{100}
{
    \zA{}
}
{
    \eA{But what she wanted to tell him is not yet known, so, reader, listen to the
        circumstances recorded in the subsequent chapter. END OF BOOK II.
        [transcriber's note: The second volume of this translation ends thus, and no
        more of it was ever published.] ERRATA [original book lists no errata; these
        were found during Project Gutenberg proofreading. The format is imitated
        from the list actually appearing at the end of volume I. If a word is split
        across a line or page then the line or page given is that on which the
        erroneous part of the word appears. On several occasions the book uses
        nested double quotes. One person, speaking, quotes another person, speaking.
        "This example," the proofreader said, "is of when my friend told me, "Don't
        take any wooden nickels."}
}
{
}

\Verse{101}
{
    \zA{}
}
{
    \eA{So I have always been careful." When these were found, the inner quotes were
        changed to single quotes for increased clarity. Such changes are not noted
        in the errata. A few other corrections to punctuation are noted below, but
        most are not. The following are not misspellings: "dumfoundered" "parricide"
        "nobble" "finicking". "shewing" was very moldy at the time this was written
        but still not deceased. The Oxford English Dictionary, Second Edition, was
        used as the authority for spellings. I don't know about "per mensem" Chapter
        XXXVI page 180, line 18. I don't know about "titify" Chapter XL page 258,
        line 21. ]  Chap.}
}
{
}

\Verse{102}
{
    \zA{}
}
{
    \eA{XXV Page 8 Line 29: doesn't \_not\_ does'nt  XXVII " 37 " 10: peccadilloes
        \_not\_ peccadiloes  XXVIII " 64 " 6: on \_not\_ ou  XXVIII " 67 " 19:
        enumeration \_not\_ enuneration  XXX " 95 " 29: them," \_not\_ them." XXX "
        100 " 24: mustn't \_not\_ musn't  XXXI " 109 " 32: needn't \_not\_ need'nt
        XXXII " 119 " 40: eh!" \_not\_ eh! XXXII " 120 " 30: "Who \_not\_ Who  XXXII
        " 128 " 13: stitch \_not\_ stich  XXXIII " 137 " 2: fidgetted \_not\_
        figetted  XXXIV " 147 " 28: promptly \_not\_ promply  XXXIV " 155 " 32:
        questions?" \_not\_ questions?}
}
{
}

\Verse{103}
{
    \zA{}
}
{
    \eA{XXXIV " 157 " 7: contrariwise \_not\_ contrarivise  XXXV " 163 " 4: eat,"
        \_not\_ eat"  XXXV " 163 " 13: successive \_not\_ succcessive  XXXV " 163 "
        35: forty \_not\_ fourty  XXXV " 171 " 12: birthday \_not\_ brithday  XXXVI
        " 180 " 2: tael. \_not\_ tael." XXXVI " 190 " 20: birthday \_not\_ brithday
        XXXVII " 194 " 18: comes \_not\_ come's  XXXVII " 198 " 10: To-morrow
        \_not\_ To-morow  XXXVII " 199 " 32: "Well," \_not\_ "Well",  XXXVII " 199 "
        33: done." \_not\_ done? XXXVII " 199 " 40: fairest \_not\_ fairiest  XXXVII
        " 206 " 13: mustn't \_not\_ musn't  XXXVII " 207 " 36: get \_not\_ ged
        XXXVII " 211 " 16: do?" \_not\_ do? XXXVIII " 219 " 6: stomachaches."
        \_not\_  stomachaches.}
}
{
}

\Verse{104}
{
    \zA{}
}
{
    \eA{XXXVIII " 228 " 13: while \_not\_ whily  XXXIX " 232 " 5: with?" \_not\_
        with?,'  XXXIX " 237 " 9: conscious \_not\_ concious  XXXIX " 242 " 1:
        temple." \_not\_ temple. XL " 245 " 38: little \_not\_ litte  XL " 248 " 11:
        silk." \_not\_ silk?" XL " 254 " 12: They're \_not\_ The're  XL " 255 " 8:
        autograph \_not\_ authograph  XL " 257 " 16: mustn't \_not\_ musn't  XL "
        258 " 13: fogies \_not\_ foggies  XL " 258 " 20: predilection \_not\_
        predeliction  XL " 258 " 35: curtains." \_not\_ curtains. XL " 258 " 39:
        enough." \_not\_ enough.}
}
{
}

\Verse{105}
{
    \zA{}
}
{
    \eA{XL " 263 " 8: peony \_not\_ peone  XLI " 278 " 11: haven't \_not\_ have'nt
        XLII " 282 " 4: haven't \_not\_ have'nt  XLII " 282 " 19: haven't \_not\_
        have'nt  XLII " 283 " 14: ensconce \_not\_ ensconse  XLII " 284 " 26:
        medicine \_not\_ medecine  XLII " 284 " 39: medicines \_not\_ medecines
        XLII " 285 " 27: medicines \_not\_ medecines  XLII " 288 " 5: aren't \_not\_
        are'nt  XLII " 290 " 27: locust \_not\_ lucust  XLII " 290 " 27: feed.'"
        \_not\_ feed.' XLIII " 309 " 31: grandiloquent \_not\_ grandeloquent  XLIV "
        314 " 12: shouldn't \_not\_ should'nt  XLIV " 316 " 4: mustn't \_not\_
        must'nt  XLIV " 317 " 6: employed the \_not\_ employed on  the  XLIV " 322 "
        3: differed \_not\_ differred  XLIV " 322 " 31: swelled \_not\_ swole  XLIV
        " 323 " 15: unhappiness \_not\_ uuhappiness  XLV " 337 " 30: ginseng \_not\_
        ginsing  XLV " 338 " 22: medicines \_not\_ medecines  XLV " 343 " 30:
        uselessly \_not\_ uselesly  XLVI " 352 " 26: mightn't \_not\_ mighn't  XLVII
        " 372 " 32: friendship \_not\_ frienship  XLVII " 378 " 3: proffered \_not\_
        proferred  XLVIII " 380 " 21: worldly \_not\_ wordly  XLVIII " 386 " 4:
        antithetical \_not\_ antetithical  XLVIII " 386 " 23: Ling \_not\_ Ling,
        XLVIII " 386 " 23: smile \_not\_ smiled  XLVIII " 386 " 35: stanzas \_not\_
        stanaas  XLVIII " 389 " 24: cockatoo \_not\_ cuckatoo  XLVIII " 391 " 27:
        'Tis \_not\_ T'is  XLVIII " 391 " 31: 'Tis \_not\_ T'is  XLIX " 393 " 34:
        would'st \_not\_ woulds't  XLIX " 393 " 37: 'tis \_not\_ t'is  XLIX " 401 "
        1: simultaneously \_not\_  simultaneouly  L " 411 " 25: 'tis \_not\_ t'is  L
        " 413 " 17: 'tis \_not\_ t'is  L " 415 " 35: But by and bye \_not\_ But and
        bye  L " 417 " 17: 'tis \_not\_ t'is  L " 417 " 17: 'tis \_not\_ 't'is  [yes
        twice in the same line]  L " 417 " 25: 'tis \_not\_ t'is  L " 418 " 10:
        haven't \_not\_ have'nt  L " 423 " 38: blossom \_not\_ blosson  LI " 437 "
        37: matter.'"}
}
{
}

\Verse{106}
{
    \zA{}
}
{
    \eA{\_not\_ matter." LII " 446 " 21: medicine \_not\_ medecine  LII " 446 " 27:
        medicines \_not\_ medecines  LII " 449 " 5: medicines \_not\_ medecines  LII
        " 460 " 3: anniversary \_not\_ anniversay  LIII " 462 " 13: perspiring
        \_not\_ prespiring  LIII " 464 " 7: peonies \_not\_ peones  LIII " 468 " 23:
        haven't \_not\_ have'nt  LIII " 471 " 39: Apparent \_not\_ Apparrent  LIII "
        476 " 9: homage \_not\_ hommage  LIII " 476 " 14: consonant \_not\_
        consonnant  LIV " 487 " 5: trod \_not\_ trode  LIV " 487 " 12: "This \_not\_
        This  LIV " 488 " 36: Isn't \_not\_ Is'nt  LIV " 490 " 15: me?" \_not\_ me?
        LIV " 490 " 19: say, \_not\_ say,"  LIV " 491 " 23: comfortable \_not\_
        confortable  LIV " 495 " 12: exhilarated \_not\_ exhilerated  LIV " 495 "
        19: smilingly \_not\_ similingly  LV " 503 " 10: and \_not\_ aud  LV " 507 "
        32: Mrs. \_not\_ "Mrs.}
}
{
}

\Verse{107}
{
    \zA{}
}
{
    \eA{LV " 507 " 33: making \_not\_ make  LVI " 525 " 27: Aren't \_not\_ Are'nt
        LVI " 529 " 18: mustn't \_not\_ musn't  LVI " 535 " 20: notwithstanding
        \_not\_  nothwithstanding  LVI " 536 " 36: aren't \_not\_ are'nt}
}
{
}
